\section{Depth: the model as its own draft}
\label{sec:depth}

The complaint addressed in this section is that speculative decoding
\cite{leviathan2023,chen2023} requires a second, separately trained network. We
show that correctness never required one (Theorem~\ref{thm:exact}), that the
acceptance rate of an early-exit draft is controlled by a measurable property of
the residual stream (Theorem~\ref{thm:accept}), and that under the cost model of
Definition~\ref{def:cost} self-drafting strictly dominates any external draft of
equal relative cost, in \emph{every} regime $w$ (Theorem~\ref{thm:dominance}).

\subsection{Correctness does not constrain the draft}

\begin{theorem}[Exactness for an arbitrary draft]
\label{thm:exact}
Let $p, q$ be any distributions on a common countable alphabet $\Vocab$, with
the convention $r = p$ when $q = p$. Then the output of a speculative step
(Definition~\ref{def:specround}) is distributed exactly as $p$, and the
acceptance probability is
\begin{equation}
  \Prob[\text{accept}] \;=\; \sum_{x} \min\{p(x), q(x)\} \;=\; 1 - \TV(p,q) .
  \label{eq:alpha}
\end{equation}
\end{theorem}

\begin{proof}
Fix $x \in \Vocab$. The token $x$ is emitted either by being drawn from $q$ and
accepted, or by being drawn from the residual \eqref{eq:residualdist} after a
rejection. The first event has probability
$q(x)\min\{1, p(x)/q(x)\} = \min\{p(x), q(x)\}$, with the convention that the
term is $0$ when $q(x) = 0$. Summing over $x$,
\begin{equation}
  \Prob[\text{accept}] = \sum_x \min\{p(x),q(x)\}
  = \sum_x \Big( p(x) - (p(x)-q(x))_+ \Big)
  = 1 - \big\|(p-q)_+\big\|_1 ,
\end{equation}
so $\Prob[\text{reject}] = \|(p-q)_+\|_1$, and the second event has probability
$\|(p-q)_+\|_1 \cdot (p(x)-q(x))_+ / \|(p-q)_+\|_1 = (p(x)-q(x))_+$. Hence
\begin{equation}
  \Prob[\text{emit } x] = \min\{p(x),q(x)\} + (p(x)-q(x))_+ = p(x),
\end{equation}
using $\min\{a,b\} + (a-b)_+ = a$. Finally
$\sum_x \min\{p,q\} = 1 - \tfrac12\|p-q\|_1$ follows from
$\min\{a,b\} = \tfrac12(a + b - |a-b|)$ and $\sum_x p = \sum_x q = 1$.
No support hypothesis is required: when $q(x) = 0 < p(x)$ the token is
unreachable by acceptance, contributes $0$ to it, and is supplied by the residual
branch with exactly the right mass $(p(x)-q(x))_+ = p(x)$.
\end{proof}

\begin{remark}[The draft is a free parameter]
\label{rem:free}
Theorem~\ref{thm:exact} places \emph{no} condition on $q$ whatsoever: $q$ may be a
smaller model, a layer-$\ell$ readout of $M$ itself
(Definition~\ref{def:logitlens}), an $n$-gram table, or noise. The output
distribution is $p$ regardless, and only the \emph{rate}
\eqref{eq:alpha} changes. Consequently no argument from correctness supports
maintaining a separate draft network; the entire design question is the ratio of
acceptance to cost, which is what Section~\ref{sec:depth-cost} formalises.
\end{remark}

\subsection{Acceptance is controlled by residual tail energy}

\begin{lemma}[Softmax is $\TV$-Lipschitz in the spread seminorm]
\label{lem:softmax}
For $z \in \R^{V}$ write $\mathrm{spr}(z) = \max_j z_j - \min_j z_j$. Then for all
$u,v \in \R^{V}$, with $\sigma = \softmax$,
\begin{equation}
  \TV\big(\sigma(u), \sigma(v)\big) \;\le\; 1 - e^{-\mathrm{spr}(u-v)}
  \;\le\; 1 - e^{-2\|u-v\|_\infty} .
\end{equation}
\end{lemma}

\begin{proof}
First suppose $\varepsilon = \|u-v\|_\infty$. For every $i$,
$\sigma(u)_i/\sigma(v)_i = e^{u_i-v_i}\sum_j e^{v_j} / \sum_j e^{u_j}$, and since
$|u_j-v_j|\le\varepsilon$ for all $j$ we get
$e^{-\varepsilon}\sum_j e^{v_j} \le \sum_j e^{u_j} \le e^{\varepsilon}\sum_j e^{v_j}$
and $e^{u_i-v_i}\in[e^{-\varepsilon},e^{\varepsilon}]$. Multiplying,
$\sigma(u)_i \ge e^{-2\varepsilon}\sigma(v)_i$, so by
$\TV(p,q) = \sum_i (q_i-p_i)_+$,
\begin{equation}
  \TV(\sigma(u),\sigma(v)) = \sum_i\big(\sigma(v)_i-\sigma(u)_i\big)_+
  \le \sum_i \sigma(v)_i\big(1-e^{-2\varepsilon}\big) = 1-e^{-2\varepsilon}.
  \label{eq:supbound}
\end{equation}
Now sharpen using shift invariance. Softmax satisfies
$\sigma(u) = \sigma(u + c\mathbf{1})$ for every $c\in\R$, so with
$z = u-v$, $s = \mathrm{spr}(z)$ and $c = -\tfrac12(\max_j z_j + \min_j z_j)$ we
have $\sigma(u) = \sigma(v + (z + c\mathbf{1}))$ and
$\|z + c\mathbf{1}\|_\infty = s/2$. Applying \eqref{eq:supbound} with
$\varepsilon = s/2$ gives $\TV \le 1 - e^{-s}$. Finally
$s \le 2\|z\|_\infty$.
\end{proof}

\begin{remark}
The sup-norm form \eqref{eq:supbound} is the one usually written down, but
$\|u-v\|_\infty$ is not shift invariant while $\TV(\sigma(u),\sigma(v))$ is, so it
is slack by up to a factor of two --- an artificial factor that would otherwise
propagate into every exponent below.
\end{remark}

\begin{lemma}[Sharp bound for normalisation]
\label{lem:sphere}
For all non-zero $a,b \in \R^{d}$,
\begin{equation}
  \Big\| \frac{a}{\|a\|} - \frac{b}{\|b\|} \Big\|
  \;\le\; \frac{2\,\|a-b\|}{\|a\| + \|b\|} .
\end{equation}
Consequently, for RMSNorm $N(h) = \sqrt d\, \Diag(g)\, h/\|h\|$,
\begin{equation}
  \big\|N(h) - N(h')\big\|
  \;\le\; \frac{2\sqrt d\,\|g\|_\infty}{\|h\| + \|h'\|}\,\big\|h-h'\big\| .
  \label{eq:normbound}
\end{equation}
\end{lemma}

\begin{proof}
Put $s = \|a\|$, $t = \|b\|$ and $c = \langle a,b\rangle/(st) \in [-1,1]$. The
squared left side is $2-2c$ and the squared right side is
$4(s^2+t^2-2stc)/(s+t)^2$. Their difference, after multiplying by $(s+t)^2>0$, is
\begin{align}
  4(s^2+t^2-2stc) - (2-2c)(s+t)^2
  &= 2s^2 + 2t^2 - 4st - 4stc + 2c s^2 + 2c t^2 \notag\\
  &= 2(s-t)^2 + 2c(s-t)^2 \;=\; 2(1+c)(s-t)^2 \;\ge\; 0,
\end{align}
since $c \ge -1$. The RMSNorm form follows from
$N(h)-N(h') = \sqrt d\,\Diag(g)\big(h/\|h\| - h'/\|h'\|\big)$ and
$\|\Diag(g)z\| \le \|g\|_\infty\|z\|$.
\end{proof}

\begin{remark}
Lemma~\ref{lem:sphere} is an exact inequality between the two endpoints and
requires no Lipschitz hypothesis. This matters: a bound of the form ``$N$ is
$\kappa$-Lipschitz on a region containing $\hs{\ell}$ and $\hs{L}$'' is vacuous
as usually stated, since a two-point set admits every constant, and repairing it
requires a convex region whose distance to the origin --- not the endpoint norms
--- sets the constant.
\end{remark}

\begin{theorem}[Acceptance lower bound]
\label{thm:accept}
Let $\alpha_\ell = 1 - \TV(p_\ell,p_L)$ be the acceptance rate of the layer-$\ell$
readout used as a draft for $M$, let $T_\ell$ be the tail energy
(Definition~\ref{def:tail}), let $N$ be RMSNorm with gain $g$, and define the
\emph{unembedding diameter}
\begin{equation}
  D_U \;=\; \max_{i,j} \big\| (W_U)_{i,:} - (W_U)_{j,:} \big\|_2
  \;\le\; 2\|W_U\|_{2,\infty} .
\end{equation}
Then
\begin{equation}
  \alpha_\ell \;\ge\;
  \exp\!\left( - \frac{2\sqrt d\,\|g\|_\infty\, D_U\, T_\ell}
                     {\|\hs{\ell}\|_2 + \|\hs{L}\|_2} \right) .
  \label{eq:acceptbound}
\end{equation}
\end{theorem}

\begin{proof}
Write $D = N(\hs{\ell}) - N(\hs{L})$ and $w_i = (W_U)_{i,:}$, so that the logit
difference has coordinates $\langle w_i, D\rangle$. By
Lemma~\ref{lem:softmax} it suffices to bound the spread:
\begin{equation}
  \mathrm{spr}\big(W_U D\big)
  = \max_{i,j} \langle w_i - w_j,\, D\rangle
  \le \max_{i,j}\|w_i - w_j\|\,\|D\|
  = D_U \|D\| ,
\end{equation}
by Cauchy--Schwarz. By Lemma~\ref{lem:sphere} applied to
$h = \hs{\ell}$, $h' = \hs{L}$ and Definition~\ref{def:tail},
$\|D\| \le 2\sqrt d\,\|g\|_\infty T_\ell / (\|\hs{\ell}\| + \|\hs{L}\|)$.
Combining and using $\alpha_\ell \ge e^{-\mathrm{spr}}$ gives the claim.
\end{proof}

\begin{remark}[Strictly stronger than the sup-norm form]
Bound \eqref{eq:acceptbound} improves the naive estimate on both factors: $D_U$
replaces $2\|W_U\|_{2,\infty}$ because only \emph{differences} of unembedding
rows affect a softmax, and $\|\hs{\ell}\|+\|\hs{L}\|$ replaces a single endpoint
norm. It is also scale-free, as it must be, since $N$ discards the scale of its
argument.
\end{remark}

\begin{remark}[What the bound is for]
\label{rem:loose}
Bound \eqref{eq:acceptbound} is monotone in the right quantity --- acceptance is
governed by how much work the model has \emph{left} to do --- but we expect it to
be numerically weak. Two inequalities give information away, and they can be
measured separately: Cauchy--Schwarz with the diameter $D_U$ assumes the worst
pair of unembedding rows aligns with the tail direction, and
Lemma~\ref{lem:softmax} is tight only when the logit gap is spread uniformly
across the vocabulary. We had expected the first to dominate;
Section~\ref{sec:experiments} reports that this expectation was wrong.
\end{remark}

\begin{proposition}[Unconditional subspace refinement]
\label{prop:rank}
Let $P$ be any orthogonal projector and $Q = I - P$. Then
\begin{equation}
  \alpha_\ell \ge \exp\big(-M_P(D)\big),
  \qquad
  M_P(D) = \max_{i,j}\Big( \|P(w_i-w_j)\|\,\|PD\| + \|Q(w_i-w_j)\|\,\|QD\| \Big),
\end{equation}
and $M_P(D) \le D_U\|D\|$ always, so this is never worse than
Theorem~\ref{thm:accept}.
\end{proposition}

\begin{proof}
Orthogonality gives
$\langle w_i-w_j, D\rangle = \langle P(w_i-w_j), PD\rangle + \langle Q(w_i-w_j), QD\rangle$,
and Cauchy--Schwarz on each term gives the bound on the spread. For the
comparison, apply Cauchy--Schwarz in $\R^{2}$ to the vectors
$(\|P(w_i-w_j)\|, \|Q(w_i-w_j)\|)$ and $(\|PD\|,\|QD\|)$, whose norms are
$\|w_i-w_j\|$ and $\|D\|$ respectively.
\end{proof}

\begin{remark}[Why the one-term version fails]
\label{rem:onetermfails}
It is tempting to drop the $Q$ term, claiming that if the tail concentrates in
$\operatorname{im}P$ then $\mathrm{spr} \le A_P\|D\|$ with
$A_P = \max_{i,j}\|P(w_i-w_j)\|$. This is \emph{false} in general. Writing
$B_P$ for the corresponding $Q$-norm and $c = \|PD\|^2/\|D\|^2$ for the captured
fraction, the one-term bound holds iff
$B_P\sqrt{1-c} \le A_P(1-\sqrt c)$; for small leakage $\eta = \sqrt{1-c}$ the
right side shrinks like $\eta^2/2$ while the left shrinks like $\eta$, so the
condition is roughly $B_P/A_P \lesssim \eta/2$. Capturing $99\%$ of the tail
energy therefore still requires the unembedding rows themselves to be $95\%$
contained in the same subspace --- a joint condition on $W_U$ and the tail, not a
statement about the tail alone. Proposition~\ref{prop:rank} avoids the issue by
keeping both terms, and is valid for a $P$ fitted to the very data it is
evaluated on, being a deterministic per-example inequality. Note also that $P$
must be built from the \emph{post-normalisation} difference $D$, not from
$\hs{L}-\hs{\ell}$, since RMSNorm's rescaling and gain change the relevant
subspace.
\end{remark}

\subsection{Cost accounting, and when self-drafting dominates}
\label{sec:depth-cost}

We first record an exact expression for round yield that makes no independence
assumption, since acceptance events at successive positions are strongly
correlated in practice.

\begin{proposition}[Exact expected yield]
\label{prop:yield}
Consider a round with lookahead $\gamma$ (Definition~\ref{def:specround}) and let
$A_j$ be the event that the first $j$ drafted tokens are all accepted,
$j = 1,\dots,\gamma$. Then the number $Y$ of tokens emitted in the round
satisfies
\begin{equation}
  Y = 1 + \sum_{j=1}^{\gamma} \mathbf{1}[A_j],
  \qquad\text{hence}\qquad
  \E[Y] = 1 + \sum_{j=1}^{\gamma} \Prob[A_j] .
  \label{eq:exactyield}
\end{equation}
\end{proposition}

\begin{proof}
Acceptance halts at the first rejection, so the number of accepted tokens is
$K = \max\{j : A_j \text{ holds}\}$ (with $K=0$ if $A_1$ fails), and $A_j$ holds
if and only if $j \le K$. Therefore $\sum_{j=1}^{\gamma}\mathbf{1}[A_j] = K$.
The round emits the $K$ accepted tokens plus exactly one further token --- drawn
from the residual \eqref{eq:residualdist} if $K < \gamma$, and from the target at
the final position if $K = \gamma$ --- so $Y = K+1$. Linearity of expectation
gives the second identity; no independence is used.
\end{proof}

\begin{corollary}[Geometric special case]
\label{cor:geom}
If acceptance events are i.i.d.\ with rate $\alpha$, then
$\Prob[A_j] = \alpha^{j}$ and
$\E[Y] = (1-\alpha^{\gamma+1})/(1-\alpha)$.
\end{corollary}

\begin{proof}
Substitute into \eqref{eq:exactyield}:
$\E[Y] = 1 + \sum_{j=1}^\gamma \alpha^j
 = 1 + \frac{\alpha - \alpha^{\gamma+1}}{1-\alpha}
 = \frac{1 - \alpha^{\gamma+1}}{1-\alpha}$.
\end{proof}

\begin{remark}
Correlated acceptance makes $\Prob[A_j] \ge \prod_{i \le j}\alpha_i$ typical ---
easy contexts accept runs, hard contexts reject early --- so
Corollary~\ref{cor:geom} evaluated at the mean acceptance rate generally
\emph{under}-estimates yield. Reported speedups computed from
Corollary~\ref{cor:geom} are therefore conservative, which is the safe direction
for the comparison below, since both schemes are scored the same way.
\end{remark}

\begin{definition}[Head fraction]
\label{def:headfrac}
Let $u \in [0,1)$ be the fraction of a full single-position forward pass
attributable to the unembedding $W_U$, so that the $L$ blocks jointly account for
$1-u$. Every draft step must evaluate $W_U$, since a draft is useless unless a
token can be sampled from it.
\end{definition}

\begin{proposition}[Round cost]
\label{prop:cost}
Let $\rho = \ell/L$ and $D_k = w + (1-w)k$. Drafting evaluates blocks $1..\ell$
at positions $t,\dots,t+\gamma-1$. Verification requires full depth at all of
$t,\dots,t+\gamma$; the final position $t+\gamma$, which supplies the bonus
token, is never touched by drafting. Two verification strategies exist:
\begin{align}
  C^{\mathrm{nr}}_{\mathrm{self}}
    &= \gamma\big[(1-u)\rho + u\big] \;+\; D_{\gamma+1},
    \label{eq:cnoreuse}\\
  C^{\mathrm{re}}_{\mathrm{self}}
    &= \gamma\big[(1-u)\rho + u\big]
     \;+\; (1-u)(1-\rho)D_{\gamma} \;+\; (1-u) \;+\; u D_{\gamma+1},
    \label{eq:creuse}\\
  C_{\mathrm{ext}}
    &= \gamma c_d \;+\; D_{\gamma+1},
    \label{eq:cext}
\end{align}
and $C_{\mathrm{self}} = \min\{C^{\mathrm{nr}}_{\mathrm{self}},
C^{\mathrm{re}}_{\mathrm{self}}\}$.
\end{proposition}

\begin{proof}
Each of the $\gamma$ draft steps evaluates blocks $1..\ell$ and the head at a
single position, costing $(1-u)\rho + u$ independently of $w$. For
\eqref{eq:cnoreuse}, verification is one batched full-depth pass over the
$\gamma+1$ positions, costing $D_{\gamma+1}$. For \eqref{eq:creuse}, cached keys
and values let the $\gamma$ drafted positions resume at block $\ell+1$, costing
$(1-u)(1-\rho)D_\gamma$, but position $t+\gamma$ has nothing cached and needs a
full-depth pass of its own, costing $(1-u)$; the head is evaluated at all
$\gamma+1$ positions. The external draft shares no blocks with $M$.
\end{proof}

\begin{remark}[Sanity check, and an error it catches]
\label{rem:sanity}
At $\rho = 1$ the draft is the target, so $\alpha = 1$ and the scheme degenerates
to ordinary autoregressive decoding, forcing $S = 1$. Both \eqref{eq:cnoreuse}
and \eqref{eq:creuse} give $C = \gamma+1$ there at $w=1$, matching the yield
$\gamma+1$. A cost model that omits the bonus position instead returns
$S = (\gamma+1)/(\gamma+u) > 1$, manufacturing speedup from nothing. An earlier
version of this paper contained exactly that error, and it was caught only by
running the experiment.
\end{remark}

\begin{proposition}[Cache reuse is regime-dependent, and useless when memory-bound]
\label{prop:reuse}
At $w = 1$, $C^{\mathrm{re}}_{\mathrm{self}} \ge C^{\mathrm{nr}}_{\mathrm{self}}$
for every $\rho \in [0,1]$, with equality only at $\rho = 1$. At $w = 0$ the
inequality reverses for every $\rho > 0$.
\end{proposition}

\begin{proof}
At $w=1$ we have $D_k = 1$ for all $k$, so
$C^{\mathrm{re}} - C^{\mathrm{nr}} = (1-u)(1-\rho) + (1-u) + u - 1 = (1-u)(1-\rho) \ge 0$.
At $w=0$, $D_k = k$, and
$C^{\mathrm{re}} - C^{\mathrm{nr}} = (1-u)\big[(1-\rho)\gamma + 1 - (\gamma+1)\big]
 = -(1-u)\rho\gamma < 0$.
\end{proof}

\begin{remark}[Retraction]
\label{rem:kvreuse}
An earlier version argued that cache reuse is the structural advantage of
self-drafting, since skipping blocks $1..\ell$ also skips streaming their
parameters. The reasoning is incomplete and the conclusion is wrong: reuse
splits verification into two passes at \emph{different depths}, and in the
memory-bound regime a pass costs a full weight stream however many positions it
covers, so the lower blocks are paid for twice. At $u=0$, $\rho=0.25$,
$\gamma=2$ reuse costs $2.25$ against $1.50$ without it. The genuine advantage of
self-drafting when memory-bound is simply that there is no second model to
stream; reuse belongs to the compute-bound case.
\end{remark}

\begin{theorem}[Self-drafting dominates]
\label{thm:dominance}
Fix $\gamma \ge 1$, $w \in [0,1]$, $\rho \in (0,1)$. If an external draft model
has relative cost $c_d \ge u + (1-u)\rho$ and acceptance
$\alpha_d \le \alpha_\ell$, then early-exit self-drafting at layer $\ell$ attains
at least as great a speedup, and strictly greater when
$c_d > u + (1-u)\rho$.
\end{theorem}

\begin{proof}
Comparing \eqref{eq:cext} with \eqref{eq:cnoreuse},
$C_{\mathrm{ext}} - C^{\mathrm{nr}}_{\mathrm{self}}
 = \gamma[c_d - u - (1-u)\rho] \ge 0$, with equality exactly when
$c_d = u + (1-u)\rho$, and $C_{\mathrm{self}} \le C^{\mathrm{nr}}_{\mathrm{self}}$.
Yield is non-decreasing in the acceptance events by
Proposition~\ref{prop:yield}, so $\E[Y\mid\alpha_\ell] \ge \E[Y\mid\alpha_d] > 0$.
\end{proof}

\begin{remark}
The strictness claimed in an earlier version came from the reuse term retracted
in Remark~\ref{rem:kvreuse}. What survives is weak dominance. Not modelled, and
further favouring the self-draft: an external draft occupies additional resident
memory and requires vocabulary alignment between the two models.
\end{remark}

\begin{corollary}[The head is a hard ceiling]
\label{cor:ceiling}
Even with perfect acceptance and $\rho \to 0$, the speedup of self-drafting
satisfies
\begin{equation}
  S \;\le\; \frac{\gamma+1}{\gamma u + D} \;\xrightarrow[w \to 1]{}\;
  \frac{\gamma+1}{\gamma u + 1} .
\end{equation}
For Qwen3-0.6B, whose tied unembedding is $151936 \times 1024$ of roughly
$5.96\times10^{8}$ parameters, $u \approx 0.26$, giving $S \le 2.45$ at
$\gamma = 4$.
\end{corollary}

\begin{remark}[Why omitting the head is not a rounding error]
\label{rem:headmatters}
Writing $C_{\mathrm{self}} = \gamma\rho + 1$, as is common, assumes a draft step
streams only the first $\ell$ blocks. It must also stream $W_U$. At
$\gamma = 4$ the omission inflates the reported speedup by a factor of $1.07$ at
$\rho = 0.75$, $1.17$ at $\rho = 0.5$ and $1.39$ at $\rho = 0.25$ --- largest
exactly where early exit is most attractive. Since our pre-registered kill
criterion for this thesis is $\max S \le 1$, an error of this size and sign can
invert the conclusion, and small models with large vocabularies are the worst
case rather than a corner case.
\end{remark}

\begin{corollary}[Shortlisted draft head]
\label{cor:shortlist}
Let $S \subseteq \Vocab$ be chosen independently of the current logits, and let
the draft be $q_S(x) = p_\ell(x)/p_\ell(S)$ for $x \in S$ and $0$ otherwise.
Then decoding remains exact, the head cost falls from $u$ to $u|S|/V$, and
\begin{equation}
  \alpha_S \;\ge\; \alpha_\ell - p_L(\Vocab \setminus S) .
\end{equation}
\end{corollary}

\begin{proof}
Exactness is Theorem~\ref{thm:exact}, which requires no support condition --- this
is where its full generality is used. For the rate, $q_S \ge p_\ell$ pointwise on
$S$, so
$\alpha_S = \sum_{x\in S}\min\{p_L,q_S\} \ge \sum_{x \in S}\min\{p_L,p_\ell\}
 = \alpha_\ell - \sum_{x \notin S}\min\{p_L,p_\ell\} \ge \alpha_\ell - p_L(S^c)$.
The cost claim is immediate, as only the rows of $W_U$ indexed by $S$ are needed.
\end{proof}

\begin{remark}
$S$ must be fixed in advance --- a static frequency shortlist, or one conditioned
on the previous token --- because selecting a top-$k$ set requires the full logit
vector and so would not save the streaming of $W_U$. Corollary~\ref{cor:shortlist}
is the natural route around Corollary~\ref{cor:ceiling}, and it trades a bounded,
\emph{measurable} acceptance loss for the head cost.
\end{remark}

\begin{remark}[The real obstruction is the training objective]
\label{rem:offdist}
Theorem~\ref{thm:dominance} is conditional on $\alpha_\ell \ge \alpha_d$, and
for an untrained early exit this is exactly what fails: $W_U$ was fit against
$N(\hs{L})$ alone, so $N(\hs{\ell})$ is off-distribution for it and
$\alpha_\ell$ is small at useful $\rho$. This is a statement about the training
objective rather than the architecture, and is what an early-exit auxiliary loss
with layer dropout addresses \cite{elhoushi2024,zhang2023draft}. The complaint of
this section is therefore best stated not as ``speculative decoding is a hack''
but as: \emph{models are trained so that their intermediate layers are not
decodable, and a second network is then introduced to supply what the first was
never asked to expose.}
\end{remark}
